\documentclass[a4paper]{deedy-resume}
\begin{document}
\lastupdated
\namesection{Jonáš}{Kulhánek}{ % Your name
\urlstyle{same}\url{https://jkulhanek.github.io} \\
\href{mailto:jonas.kulhanek@live.com}{jonas.kulhanek@live.com} | +420 728 042 412
}
% TODO: remove
\vspace{-0.3cm}

\begin{minipage}[t]{0.33\textwidth} % The left column takes up 33% of the text width of the page

\section{Education}
\subsectionlong{Charles University}\\
\descript{MSc in Artificial Intelligence}
\location{Sep 2019 \textendash\, Jun 2021}
Faculty of Mathematics and Physics\\
Summa cum laude\\
Cum. GPA: 4.00 / 4.00\\
Major GPA: 4.00 / 4.00\\
Final exam grade: A
\sectionspace

\subsectionlong{Czech Technical University}\\
\descript{BSc in Computer Science}
\location{Jun 2016 \textendash\, Jun 2019}
Faculty of Electrical Engineering\\
Summa cum laude\\
Cum. GPA: 3.64 / 4.00\\
Major GPA: 3.74 / 4.00\\
Final exam grade: A
\sectionspace

\subsectionlong{Sungkyunkwan University}\\

\location{Sep 2017 \textendash\, Dec 2017}
College of Sciences and Engineering\\
One semester visiting student\\
Cum. GPA: 4.00 / 4.00
\sectionspace


\section{Links} 

GitHub:// \href{https://github.com/jkulhanek}{\bf jkulhanek} \\
LinkedIn:// \href{https://www.linkedin.com/in/jonaskulhanek}{\bf jonaskulhanek} \\
Google Scholar:// \href{https://scholar.google.com/citations?user=YDNzfN4AAAAJ}{\bf Jonáš Kulhánek} \\

\sectionspace

\section{Teaching}
\subsectionlong{Introduction to Artificial Intelligence}\\
\location{Teaching assistant | Summer 2020}
zero-sum games, MCTS, A*, CSP, MDP, EA, planning, knowledge representation\\

\textit{FEL, Czech Technical University}
\sectionspace


\section{Projects}
\subsectionlong{Deep RL PyTorch}\\
\location{DRL, PyTorch}
Library for training DRL agents\\
\href{https://github.com/jkulhanek/deep-rl-pytorch}{\textit{github.com/jkulhanek/deep-rl-pytorch}}
\sectionspace

\subsectionlong{LemmaTag}\\
\location{NLP, TensorFlow 2}
Implementation of SoTA lemmatizer and tagger achieving 98.75\% and 96.67\% accuracies respectively on UD Treebank\\
\href{https://github.com/jkulhanek/lemmatag}{\textit{github.com/jkulhanek/lemmatag}}
\sectionspace

\subsectionlong{DMHouse}\\
\location{DRL, Bazel}
An indoor 3D environment simulator for pre-training VN agents\\
\href{https://github.com/jkulhanek/dmhouse}{\textit{github.com/jkulhanek/dmhouse}}
\sectionspace



\vspace{-1cm}

\end{minipage} % The end of the left column
\hfill
\begin{minipage}[t]{0.66\textwidth}

\section{Selected Experience}
\runsubsection{CIIRC, Czech Technical University in Prague}
\descript{| PhD student}
\location{Sep 2021 – Present}
\vspace{\topsep}
\begin{tightitemize}
\item Supervisor \textbf{Torsten Sattler} 
\item Using priors in neural rendering and neural scene representation
\item \textbf{TensorFlow 2}, PyTorch, \ldots
\end{tightitemize}
\sectionspace

\runsubsection{CIIRC, Czech Technical University in Prague}
\descript{| Junior researcher}
\location{Aug 2020 – Aug 2021}

\begin{tightitemize}
\item Visual navigation using deep reinforcement learning (DRL)
\item Implementing SoTA DRL algorithms including \textbf{Rainbow, PAAC, AlphaZero} 
\item Publishing an open-source 3D environment simulator to train DRL agents
\item Working with \textbf{Tomas Mikolov} on cellular automata applied to NLP
\item Researching dialogue systems with pre-trained LMs (GPT2, BERT)
\item Running large distributed GPU training
\item Achieving \textbf{3rd place} in the DSTC9 end-to-end multi-domain dialogue task
\item \textbf{PyTorch}, Slurm, Numpy, ROS, \ldots
\end{tightitemize}
\sectionspace

\runsubsection{Delft University of Technology}
\descript{| Research intern}
\location{Feb 2019 – Apr 2019}

\begin{tightitemize}
\item Visual navigation using deep reinforcement learning (\textbf{DRL})
\item Working under Prof. dr. \textbf{Robert Babuska} at 3me
\item Publishing open-source DRL PyTorch library
\item Evaluating navigation algorithms on real \textbf{mobile robots} 
\item \textbf{PyTorch}, TensorFlow, ROS, Numpy, \ldots
\end{tightitemize}
\sectionspace

\runsubsection{Center for Machine Perception}
\descript{| Research intern}
\location{Feb 2018 – Jun 2018}

\begin{tightitemize}
\item Building object detectors using fully convolutional neural networks
\item Implementing RetinaNet, FasterRCNN, MaskNet (SoTA at that time)
\item Preparing object detection datasets and data pipelines
\item \textbf{TensorFlow}, Numpy, OpenCV, \ldots
\end{tightitemize}
\sectionspace

complete list at \urlstyle{same}\url{https://jkulhanek.github.io/}.

\section{Publications}
\subsectionlong{ViewFormer: NeRF-free Neural Rendering from Few Images Using Transformers}\\
\location{Mar 2022}
Jonáš Kulhánek and Erik Derner and Torsten Sattler and Robert Babuška,\\
Model capable of both novel view synthesis and camera pose estimation.\\

\textit{accepted to ECCV 2022}
\sectionspace

\subsectionlong{AuGPT: Auxiliary Tasks and Data Augmentation for End-To-End Dialogue with Pre-Trained Language Models}\\
\location{Feb 2021}
Jonáš Kulhánek and Vojtěch Hudeček and Tomáš Nekvinda and Ondřej Dušek,\\
Using pre-trained GPT2 with extensive data augmentation.\\

\textit{published in proceedings of NLP4ConvAI, EMNLP 2021}
\sectionspace

\subsectionlong{Visual Navigation in Real-World Indoor Environments Using End-to-End Deep Reinforcement Learning}\\
\location{Apr 2021}
Jonáš Kulhánek and Erik Derner and Robert Babuška,\\
Designing DRL agent with auxiliary tasks for real-world navigation. Transfering policy pretrained on custom 3D simulator to the real world.\\

\textit{published in RA-L 2021}
\sectionspace



\vspace{-1cm}
\end{minipage}

%----------------------------------------------------------------------------------------
%	SECOND PAGE (EXAMPLE)
%----------------------------------------------------------------------------------------

%\newpage % Start a new page

%\begin{minipage}[t]{0.33\textwidth} % The left column takes up 33% of the text width of the page

%\section{Example Section}

%\end{minipage} % The end of the left column
%\hfill
%\begin{minipage}[t]{0.66\textwidth} % The right column takes up 66% of the text width of the page

%\section{Example Section 2}

%\end{minipage} % The end of the right column

%----------------------------------------------------------------------------------------

\end{document}
